\section{Variants}\label{sec:variants}

This section records what the same protocol certifies for other degree notions, what one run
costs when several constraints are imposed at once, and what separate runs give instead.

\subsection{Individual and weighted degrees}\label{sec:individual}

The weighted degree covered by Theorem~\ref{thm:main} includes the two degrees usually asked for,
since $w=(1,\dots,1)$ gives the total degree and $w=e_l$ the degree in $x_l$. For the second
choice the protocol certifies an individual degree bound of its own.

\begin{remark}[Individual degrees]
Since $C_{w,D}\subseteq\RS[\F,L_0,k^m]$, acceptance also certifies
$\deg_{x_l}<k$ for all $l$; taking $w=e_l$ and $D=b<k$ certifies $\deg_{x_l}f\le b$.
\end{remark}

The bound $b$ is free below $k$: the same commitment and the same protocol certify
$\deg_{x_l}f\le b$ for every $b<k$, and not only for $b=k-1$. The caps of
Construction~\ref{def:protocol} read $d_i=0$ for $i<l$ and $d_i=b$ for $i\ge l$ in this case, so the
words before level $l$ are scalar and the constraint costs no $T$-degrees. A bound $b\ge k-1$ adds nothing, because the commitment already restricts $f$ to the box
$\deg_{x_l}<k$ for all $l$: the range of individual degrees is fixed by the encoding, and the
test certifies bounds inside it.

\subsection{Several constraints in one run}\label{sec:multi}

Let $p\ge1$, let $w^{(1)},\dots,w^{(p)}\in\N^m$ be weights and $D_1,\dots,D_p$ bounds, and
write $\mathbf w:=(w^{(1)},\dots,w^{(p)})$, $\mathbf D:=(D_1,\dots,D_p)$, and
$d^{(s)}_i:=\min\big(D_s,(k-1)\sum_{l\le i}w^{(s)}_l\big)$, the cap of \cref{def:protocol}
for the $s$-th constraint. The target code is
$C_{\mathbf w,\mathbf D}:=\bigcap_sC_{w^{(s)},D_s}$. One run of the protocol with $p$ formal
variables proves all $p$ constraints at once.

{\sloppy
\begin{remark}[One run with $p$ formal variables]\label{rem:multi}
Replace $\F[T]$ by $\F[T_1,\dots,T_p]$, the challenge $r_iT^{w_i}$ by
$r_i\prod_sT_s^{\,w^{(s)}_i}$, the graded degree by the vector of degrees
$\big(a_s+\sum_{l>i}w^{(s)}_l\alpha_l\big)_s$ of a monomial
$\prod_sT_s^{a_s}x^\alpha$, the condition $\gdeg\le D$ by ``$\le D_s$ in every
coordinate'', the formats $\F[T]_{\le d_i}$ by $\prod_s\F[T_s]_{\le d^{(s)}_i}$, and the final
check $\deg g\le D$ by $\deg_{T_s}g\le D_s$ for all $s$. Then \cref{lem:inj},
\cref{lem:fold-sub}, \cref{lem:stab}, \cref{lem:unfold}, \cref{lem:round}, Theorem~\ref{thm:main},
Corollary~\ref{cor:ud} and Proposition~\ref{prop:complete}(a)(b) hold verbatim for
$C_{\mathbf w,\mathbf D}$, with $t:=\prod_s\big(1+\max(D_s,\,d^{(s)}_{i-1}+w^{(s)}_i(k-1))\big)$
in item~4 of \cref{lem:round}.

The proofs use only three facts about $\F[T]$: it is an integral domain (cancellation in
\cref{lem:unfold} and \cref{lem:round}); divisibility by a monomial is read off coefficientwise
(\cref{lem:div} and \cref{lem:round}, item~4); and the value spaces are finite-dimensional over
$\F$, which is all \cref{lem:ext} needs. All three hold for $\F[T_1,\dots,T_p]$ with bounded
degrees.
\end{remark}\par}
% = R1

The price is paid in the size of the messages: a value of $u_i$ is now an element of
$\prod_s\F[T_s]_{\le d^{(s)}_i}$, that is, a vector of $\prod_s(d^{(s)}_i+1)$ field elements,
and the analysis of \cref{lem:round} runs with the bound $t$ above in place of
$1+\max(D,d_{i-1}+w_i(k-1))$. \Cref{sec:costs} compares this with separate runs on a
concrete instance.

\subsection{Separate runs}\label{sec:separate}

\begin{remark}[Separate runs on one commitment]\label{rem:separate}
Running $\Pi_{w^{(s)},D_s}$ separately on the same $u_0$ costs the sum of the $p$ runs
instead of the product $t$. Below the unique-decoding radius it certifies the same
statement: if $\agree(u_0,C_{w^{(s)},D_s})>\frac{1+\rho}2$ for all $s$, the codewords
achieving these agreements are all the unique closest codeword of
$\RS[\F,L_0,k^m]$, which therefore lies in $C_{\mathbf w,\mathbf D}$. Above that radius,
separate runs only certify closeness to each $C_{w^{(s)},D_s}$, possibly through different
codewords of the list; \cref{rem:multi} is then the right tool.
\end{remark}
% = rem:separate

\subsection{Costs on an instance}\label{sec:costs}

We take the instance of \cref{sec:numbers}: $N=2^{20}$, $k=4$, $m=8$, so $|L_i|=2^{20-2i}$,
and two constraints, $\deg_{\mathrm{tot}}f\le6$ ($w^{(1)}=(1,\dots,1)$, $D_1=6$) and
$\deg_{x_1}f\le2$ ($w^{(2)}=e_1$, $D_2=2$). The caps are $d^{(1)}=(3,6,6,\dots,6)$ and
$d^{(2)}=(2,2,\dots,2)$ at levels $1,\dots,8$. \Cref{tab:costs} counts field elements; the
committed word $u_0$, common to all columns, is not counted in the proof length.

\begin{table}[ht]
\centering\small
\renewcommand{\arraystretch}{1.2}
\begin{tabular}{@{}p{4.2cm}p{2.9cm}p{3.2cm}p{2.4cm}@{}}
\hline
 & one run, $T_1,T_2$ & two separate runs & FRI only\\
\hline
value size at levels $1$; $\ge2$ & $12$; $21$ & $4+3$; $7+3$ & $1$; $1$\\
proof length (oracles $u_1,\dots,u_7$) & $4\,980\,288$ & $1\,660\,096+1\,048\,512$ & $349\,504$\\
last message & $21$ & $7+3$ & $1$\\
reads per repetition & $556$ & $188+88=276$ & $32$\\
certifies, up to Johnson & both constraints, one codeword & each constraint separately & $\deg_{x_l}<4$\\
certifies, below unique decoding & both constraints & both constraints & $\deg_{x_l}<4$\\
\hline
\end{tabular}
\caption{Costs of testing $\deg_{\mathrm{tot}}f\le6$ and $\deg_{x_1}f\le2$ at $N=2^{20}$,
$k=4$, $m=8$. Reads per repetition are $k\sum_{i=1}^m s_{i-1}$, where $s_i$ is the value size
at level $i$ and $s_0=1$.}
\label{tab:costs}
\end{table}

One run costs about twice as much as the two separate runs here, because the value sizes
multiply instead of adding. The gap grows with the number of constraints. Bounding every
individual degree by $b=2$ needs $m=8$ weights $e_1,\dots,e_8$: in one run the value size at
level $i$ is $3^i$, up to $6561$ at the last level, and a repetition reads
$4\sum_{i=1}^{8}3^{i-1}=13\,120$ field elements; eight separate runs read
$\sum_{l=1}^{8}4\,(l+3(8-l))=480$ field elements in total.

The choice therefore depends on the proximity regime. Below the unique-decoding radius,
separate runs certify the intersection (\cref{rem:separate}) and are cheaper. Beyond it, only
one run with several formal variables certifies that a single polynomial satisfies all the
constraints, at the cost shown in \cref{tab:costs}.
